\documentclass[11pt]{article}

\usepackage[margin=1in]{geometry}
\usepackage{booktabs}
\usepackage{array}
\usepackage{longtable}
\usepackage{tabularx}
\usepackage{enumitem}
\usepackage{amsmath}
\usepackage{siunitx}
\usepackage{xcolor}
\usepackage{hyperref}
\usepackage{caption}

\hypersetup{
  colorlinks=true,
  linkcolor=blue!55!black,
  citecolor=blue!55!black,
  urlcolor=blue!55!black
}

\setlength{\parskip}{0.55em}
\setlength{\parindent}{0pt}
\setlist[itemize]{topsep=0.25em,itemsep=0.2em,leftmargin=1.5em}
\setlist[enumerate]{topsep=0.25em,itemsep=0.2em,leftmargin=1.5em}

\newcommand{\code}[1]{\texttt{#1}}
\newcommand{\robust}{\textcolor{green!45!black}{\textbf{robust}}}
\newcommand{\unstable}{\textcolor{orange!70!black}{\textbf{unstable}}}
\newcommand{\notrobust}{\textcolor{gray!70!black}{not robust}}

\title{\textbf{Visual Encoder Pretraining Leaves a Fingerprint on Generated Captions}\\
\large A controlled frozen-decoder captioning probe}
\author{TTIC Embeddings Project}
\date{May 24, 2026}

\begin{document}
\maketitle

\begin{abstract}
This report summarizes a controlled captioning experiment comparing four frozen
visual encoders: CLIP, SigLIP, DINOv2, and MAE. The central question is whether
the pretraining objective and data of a visual encoder affect the linguistic
form of captions produced by a fixed language model. To isolate the encoder, the
experiment freezes both the visual encoder and GPT-2 medium, trains only a small
prefix adaptor, and keeps the COCO data, decoding procedure, and evaluation
metrics fixed across conditions. The final analysis aggregates three generation
seeds to distinguish stable encoder effects from seed-specific artifacts. The
cleanest comparable-quality comparison excludes MAE because it repeatedly fails
the CIDEr and SPICE quality preconditions. In the CLIP+SigLIP versus DINOv2 contrast,
projective spatial language is robust across beam and nucleus decoding, while
caption length is robust under beam search only. In both cases the direction is
the same and counterintuitive: the self-supervised encoder (DINOv2) produces
\emph{more} projective spatial language and \emph{longer} captions than the
language-supervised encoders. These differences are statistically robust ---
consistent in sign and above the pre-registered effect-size floor in every
seed --- but small in absolute terms. Pairwise contrasts show CLIP carries
most of the pooled effect: SigLIP sits between CLIP and DINOv2, so the headline
should be read as ``CLIP, and to a lesser extent SigLIP, differ from DINOv2,''
not as a clean category-level claim. Effects for adjective density and
head-noun specificity appear in one seed but are not stable across seeds.
\end{abstract}

\section{Research Question}

The experiment asks:

\begin{quote}
If the language model, adaptor architecture, training data, and decoding
procedure are fixed, how does visual encoder pretraining change generated
captions?
\end{quote}

The project is not intended to build a state-of-the-art captioner. Its value is
as a controlled probe. Most vision-language systems vary several components at
once: visual encoder, connector, LLM, instruction data, training recipe, and
decoding procedure. Here, the visual encoder is the main experimental variable.
If captions differ under this bottleneck, the difference is evidence about what
the frozen visual representation makes easier or harder for the same frozen
decoder to say.

\section{Experimental Design}

\subsection{System Architecture}

Each condition uses the same high-level architecture:

\[
\text{image} \rightarrow \text{frozen visual encoder} \rightarrow
\text{trainable prefix adaptor} \rightarrow \text{frozen GPT-2} \rightarrow
\text{caption}.
\]

Only the prefix adaptor is trained. The visual encoders and GPT-2 are frozen.
The adaptor maps visual patch tokens into \(k=10\) soft prompt vectors in the
GPT-2 embedding space. This makes the experiment closer to a controlled
adapter-based multimodal probe than to a fully optimized BLIP-2 or LLaVA-style
system.

\subsection{Encoders}

\begin{table}[h]
\centering
\caption{Frozen visual encoders used in the experiment.}
\label{tab:encoders}
\begin{tabularx}{\textwidth}{@{}l l l X@{}}
\toprule
Encoder & Checkpoint & Supervision & Pretraining intuition \\
\midrule
CLIP & \code{openai/clip-vit-large-patch14} & image-text &
Contrastive image-text pretraining; representations are explicitly aligned
with natural-language concepts. \\
SigLIP & \code{google/siglip-large-patch16-256} & image-text &
Image-text pretraining with a sigmoid pairwise loss; also language-aligned,
but with a different loss and data recipe from CLIP. \\
DINOv2 & \code{facebook/dinov2-large} & image-only &
Self-supervised visual consistency / self-distillation; learns semantic visual
structure without captions. \\
MAE & \code{facebook/vit-mae-large} & image-only &
Masked patch reconstruction; learns visual structure from recovering missing
image patches, less directly aligned with caption semantics. \\
\bottomrule
\end{tabularx}
\end{table}

\subsection{Data and Outputs}

The adaptor training uses COCO \code{train2017}. Caption generation and
evaluation use COCO \code{val2017}. Each seed generates captions for:

\[
4 \text{ encoders} \times 2 \text{ decoders} \times 5000 \text{ images}
= 40{,}000 \text{ captions}.
\]

The experiment generates, scores, and analyzes three seeds, then aggregates
their statistical verdicts to separate robust findings from seed-sensitive
ones. The main result files are:

\begin{itemize}
  \item \code{captions/captions\_seed\{1,2,3\}.jsonl}: generated captions.
  \item \code{captions/scores\_seed\{1,2,3\}.csv}: per-caption metric scores.
  \item \code{captions/analysis\_seed*\_*.csv}: per-seed statistical tests.
  \item \code{captions/analysis\_aggregate\_*.csv}: cross-seed summaries.
  \item \code{captions/quality\_summary.txt}: CIDEr quality precondition summary.
\end{itemize}

\subsection{Training Protocol}

A 5000-image holdout drawn from \code{train2017} with a fixed
\code{holdout\_seed=0} (shared across all encoder $\times$ seed runs) supplies
the validation set used for early-stopping val PPL. Sharing one holdout across
conditions removes a fairness confound: every encoder is evaluated against the
same images during training. Each adaptor trains with AdamW
(lr $1\!\times\!10^{-4}$, weight decay 0.01), batch size 256, 500 warmup steps,
and a cosine schedule to a 35{,}000-step cap, with early-stopping patience of 5
validation evaluations (every 1000 steps). Generation uses beam search with
width 5 and nucleus sampling with $p=0.9$, both capped at
\code{max\_new\_tokens=30} (see \code{configs/base.yaml}).

\section{Metrics}

The caption-style analysis scores every generated caption on several linguistic
or semantic metrics:

\begin{itemize}
  \item \textbf{Caption length}: token length.
  \item \textbf{Projective spatial density}: density of spatial expressions
  such as ``left of'', ``behind'', ``next to'', ``in front of'', ``above'', and
  ``below''.
  \item \textbf{Topological spatial density}: density of expressions such as
  ``in'', ``on'', ``at'', ``inside'', ``outside'', and ``with''.
  \item \textbf{Adjectives per noun}: adjective density normalized by nouns.
  \item \textbf{Head noun minimum depth}: a rough abstraction/specificity
  measure for the caption head noun.
  \item \textbf{Scene-object ratio}: scene words relative to object words.
  \item \textbf{VG attribute precision/recall}: overlap with Visual Genome
  attribute annotations where available.
\end{itemize}

Reference-based caption quality is checked using two complementary metrics:
CIDEr (consensus-based n-gram similarity to the COCO references) and SPICE
(semantic scene-graph overlap with the references). Both are reported with
their standard implementations: CIDEr uses the Stanford PTBTokenizer; SPICE
uses Stanford CoreNLP's scene-graph parser. SPICE 1.0 depends on the Nashorn
JavaScript engine, which was removed from the JDK in version 15, so the
scoring run is executed under JDK 11; this is recorded in the reproducibility
table. The quality precondition is a \emph{relative} cross-encoder spread
test on each reference metric: the spread must be at most 10\% of the best
encoder. Both CIDEr and SPICE are reported as preconditions, and the
headline analysis only proceeds for the subset of encoders that pass both.

\section{Statistical Analysis}

\subsection{Primary Family}

The primary test compares language-supervised encoders against
self-supervised encoders on each metric. The full four-encoder contrast is:

\[
(\text{CLIP}, \text{SigLIP}) \quad \text{vs.} \quad
(\text{DINOv2}, \text{MAE}).
\]

Because MAE fails the quality precondition, the clean headline contrast is:

\[
(\text{CLIP}, \text{SigLIP}) \quad \text{vs.} \quad \text{DINOv2}.
\]

For each metric, the analysis computes a paired nonparametric contrast over
images and reports an effect size \(r\). A result is labeled
\code{claimed\_meaningful} only when it is statistically significant after
Benjamini-Hochberg correction and meets the pre-registered effect-size floor
\(r \ge 0.10\). Following \code{methods.md}, BH-FDR is applied separately
within each family (primary vs.\ secondary) and within each
(seed, decoder) cell at $\alpha=0.05$; no further correction is applied
across seeds or decoders. The cross-seed aggregation (\S\ref{sec:agg}) is a
categorical ``all three seeds meaningful'' rule applied on top of the
per-seed BH-corrected verdicts, not an additional statistical test.

\subsection{Secondary Family}

The secondary family runs all pairwise encoder contrasts, such as CLIP versus
SigLIP and CLIP versus DINOv2. This is diagnostic: it checks whether a pooled
category effect is driven by one encoder or whether there are large
within-category differences.

\subsection{Image-Controlled Direction Check}
\label{sec:mixed-effects}

The mixed-effects model serves a role the rank-based primary family cannot: the
effect size \(r\) is unsigned, so this model is what supplies the
\emph{direction} of each effect. It also asks whether the encoder-group effect
remains after accounting for per-image baselines:

\[
\text{score} \sim \text{supervision category} + (1 \mid \text{image}).
\]

The image random intercept absorbs the fact that some images naturally produce
longer, more spatial, or more specific captions than others.

\textbf{Deviation from pre-registration.} \code{methods.md} pre-registers a
crossed-random-effects specification, $(1 \mid \text{encoder}) + (1 \mid \text{image})$.
The implementation drops the encoder random intercept for two reasons noted
in \code{methods.md}: supervision category is a deterministic function of
encoder identity (CLIP, SigLIP $\to$ language; DINOv2, MAE $\to$ self), so an
encoder random intercept is collinear with the supervision fixed effect and
the model is formally unidentifiable; and \code{statsmodels} additionally
does not fit crossed random effects, so even an identifiable variant would
not be straightforward in this stack. When even the image-only
random intercept is singular (typically because the per-image differences are
heavily zero-inflated), the code falls back to a paired per-image \(t\)-test on
language-supervised minus DINOv2 means. The point estimate equals the per-image
mean difference and is unchanged by the fallback, but the inferential structure
is not the random-effects model the pre-registration described. The Method
column in the mixed-effects table records which path was taken for each metric.
For the two robust headline metrics this matters: \code{caption\_length} falls
back in all three seeds and \code{projective\_density} falls back in two of
three.

\subsection{Cross-Seed Aggregation}
\label{sec:agg}

The final analysis aggregates seeds 1--3. For each row, the aggregate table
reports:

\begin{itemize}
  \item \(r_{\text{seed1}}, r_{\text{seed2}}, r_{\text{seed3}}\): per-seed
  effect sizes.
  \item \(\overline{r}\): mean effect size across seeds.
  \item \(\mathrm{SD}(r)\): seed variability.
  \item \code{n\_meaningful\_seeds}: number of seeds where the metric was
  meaningful.
  \item \code{claimed\_robust}: true only if all three seeds were meaningful.
  \item \code{verdict\_unstable}: true if some seeds were meaningful but not
  all.
\end{itemize}

\section{Quality Precondition and MAE}

The project requires the encoders to produce captions of roughly comparable
reference quality before claiming stylistic differences. The threshold is a
maximum beam spread of 10\% of the best encoder, applied to both reference
metrics (CIDEr and SPICE). MAE fails this gate on both metrics in all three
seeds; the language-supervised pair plus DINOv2 passes both gates in every
seed.

\begin{table}[h]
\centering
\caption{Beam CIDEr by seed (Stanford PTBTokenizer, JDK 11). The
precondition fails in every seed because MAE is lower than the other
encoders. ``Lang+SS$_{-\text{MAE}}$'' is the MAE-excluded triple (CLIP,
SigLIP, DINOv2); its spread is well under 10\% in every seed, so the
headline analysis is run on this triple.}
\label{tab:cider}
\begin{tabular}{@{}lrrrrrr@{}}
\toprule
Seed & CLIP & SigLIP & DINOv2 & MAE & Full spread & Lang+SS$_{-\text{MAE}}$ spread \\
\midrule
1 & 1.02 & 1.11 & 1.10 & 0.88 & 0.205 & 0.086 \\
2 & 1.06 & 1.13 & 1.11 & 0.89 & 0.213 & 0.060 \\
3 & 1.03 & 1.14 & 1.11 & 0.87 & 0.230 & 0.090 \\
\bottomrule
\end{tabular}
\end{table}

\begin{table}[h]
\centering
\caption{Beam SPICE by seed (Stanford CoreNLP scene-graph parser, JDK 11).
The encoder ranking is identical to CIDEr (SigLIP $\approx$ DINOv2 $>$ CLIP
$>$ MAE), and the precondition also fails for the four-encoder set in every
seed because of MAE. With MAE excluded, the SPICE spread of the remaining
triple is under 10\% in every seed --- so the MAE-exclusion decision is
supported by both reference metrics.}
\label{tab:spice}
\begin{tabular}{@{}lrrrrrr@{}}
\toprule
Seed & CLIP & SigLIP & DINOv2 & MAE & Full spread & Lang+SS$_{-\text{MAE}}$ spread \\
\midrule
1 & 0.189 & 0.201 & 0.201 & 0.169 & 0.162 & 0.062 \\
2 & 0.193 & 0.205 & 0.202 & 0.170 & 0.168 & 0.056 \\
3 & 0.192 & 0.205 & 0.200 & 0.168 & 0.180 & 0.066 \\
\bottomrule
\end{tabular}
\end{table}

This does not imply that MAE is a bad vision model in general. It means MAE is
behind in this specific frozen-encoder, small-adaptor, frozen-GPT-2 captioning
setup. MAE remains useful as a secondary finding: reconstruction-pretrained
features are less plug-and-play for caption generation than CLIP, SigLIP, and
DINOv2 under this adapter-only bottleneck. However, MAE should not be part of
the main comparable-quality style claim.

\section{Main Results}

\subsection{Beam Decoding, MAE Excluded}

The cleanest headline table is the aggregate beam primary analysis with MAE
excluded. Two metrics are robust across all three seeds: projective spatial
density and caption length. Two effects that appear in only one seed---adjectives
per noun and head noun depth---are unstable.

\begin{table}[h]
\centering
\caption{Primary aggregate results for beam decoding, MAE excluded. The final
column gives the sign of the language-supervised minus DINOv2 difference (from
the mixed-effects table); ``lower'' means the language-supervised encoders
score below DINOv2. ``Eff.\ $n$'' is the mean across seeds of the Wilcoxon
effective sample size --- the number of image pairs with non-zero rank
difference --- not the 5000 paired images. See \S\ref{sec:magnitudes}.}
\label{tab:beam-primary}
\footnotesize
\begin{tabular}{@{}lrrrrr r l l@{}}
\toprule
Metric & \(r_1\) & \(r_2\) & \(r_3\) & Mean \(r\) & SD & Eff.\ $n$ & Verdict & Lang.\ vs DINOv2 \\
\midrule
Projective spatial density & 0.276 & 0.256 & 0.345 & 0.292 & 0.047 & 1105 & \robust & lower \\
Caption length & 0.191 & 0.287 & 0.128 & 0.202 & 0.080 & 4075 & \robust & lower \\
Adjectives per noun & 0.101 & 0.058 & 0.066 & 0.075 & 0.023 & 1537 & \unstable & lower \\
Head noun min depth & 0.102 & 0.011 & 0.055 & 0.056 & 0.046 & 2154 & \unstable & lower \\
Topological density & 0.054 & 0.047 & 0.077 & 0.059 & 0.016 & 4307 & \notrobust & higher \\
Scene-object ratio & 0.009 & 0.011 & 0.028 & 0.016 & 0.010 & 1601 & \notrobust & higher \\
VG attr precision & 0.260 & 0.014 & 0.088 & 0.121 & 0.126 & 42 & \notrobust & higher \\
VG attr recall & 0.109 & 0.014 & 0.032 & 0.052 & 0.050 & 349 & \notrobust & lower \\
\bottomrule
\end{tabular}
\end{table}

The strongest beam finding is projective spatial density
(\(\overline{r}=0.292\)), meaningful in every seed: DINOv2 captions carry more
projective spatial language than the CLIP+SigLIP average. Caption length is
also robust under beam search (\(\overline{r}=0.202\)), again with DINOv2
higher. The seed aggregation narrows the story: one-seed indications for
adjective density and noun specificity are not stable enough to headline.
The pooled framing is also a partial summary: \S\ref{sec:within} below shows
that CLIP and SigLIP differ from each other on some metrics and from DINOv2
at different magnitudes, so the headline is best read as a CLIP-anchored
contrast that SigLIP echoes more weakly.

\subsection{Nucleus Decoding, MAE Excluded}

Nucleus decoding gives a narrower result. Projective spatial density remains
robust; caption length does not.

\begin{table}[h]
\centering
\caption{Primary aggregate results for nucleus decoding, MAE excluded.
``Lang.\ vs DINOv2'' gives the sign of the language-supervised minus DINOv2
difference; only the projective-density row is robust. ``Eff.\ $n$'' is the
mean Wilcoxon effective sample size across seeds.}
\label{tab:nucleus-primary}
\footnotesize
\begin{tabular}{@{}lrrrrr r l l@{}}
\toprule
Metric & \(r_1\) & \(r_2\) & \(r_3\) & Mean \(r\) & SD & Eff.\ $n$ & Verdict & Lang.\ vs DINOv2 \\
\midrule
Projective spatial density & 0.131 & 0.175 & 0.175 & 0.160 & 0.026 & 1604 & \robust & lower \\
Caption length & 0.057 & 0.005 & 0.051 & 0.038 & 0.028 & 4691 & \notrobust & higher \\
Head noun min depth & 0.067 & 0.026 & 0.013 & 0.035 & 0.028 & 3610 & \notrobust & lower \\
Scene-object ratio & 0.032 & 0.036 & 0.006 & 0.024 & 0.016 & 1883 & \notrobust & higher \\
Topological density & 0.016 & 0.005 & 0.047 & 0.023 & 0.022 & 4909 & \notrobust & higher \\
VG attr precision & 0.033 & 0.018 & 0.048 & 0.033 & 0.015 & 169 & \notrobust & higher \\
VG attr recall & 0.037 & 0.010 & 0.030 & 0.026 & 0.014 & 782 & \notrobust & \(\approx\)0 \\
Adjectives per noun & 0.005 & 0.012 & 0.018 & 0.012 & 0.006 & 3602 & \notrobust & lower \\
\bottomrule
\end{tabular}
\end{table}

This decoder interaction is useful. It suggests that projective spatial
language is the most stable encoder-induced linguistic signature. Caption
length appears more tied to beam search's stronger mode-seeking behavior.

\subsection{Image-Controlled Direction and Magnitude}

Because the rank-based effect sizes in the tables above are unsigned, the
mixed-effects table below is what establishes the \emph{direction} of every
effect reported in this paper. It also gives an image-controlled magnitude in
each metric's native unit. In the MAE-excluded beam contrast, every robust
effect points the same way: language-supervised encoders score \emph{lower}
than DINOv2 on projective spatial density and on caption length, after
controlling for image identity.

\begin{table}[h]
\centering
\caption{Image-controlled aggregate coefficients for beam decoding, MAE excluded.
Negative means language-supervised lower than DINOv2. The Method column counts,
across the three seeds, how many fit the pre-registered image random-intercept
model (\code{mixedlm}) and how many fell back to a paired per-image $t$-test
(\code{paired-t}) after the random-effects covariance went singular.}
\label{tab:mixed-effects}
\begin{tabular}{@{}lrrl@{}}
\toprule
Metric & Mean coefficient & SD & Method (seeds 1--3) \\
\midrule
Caption length & $-0.3207$ & $0.1395$ & paired-t (3/3) \\
Projective spatial density & $-0.0041$ & $0.0010$ & paired-t (2/3); mixedlm (1/3) \\
Adjectives per noun & $-0.0083$ & $0.0015$ & mixedlm (3/3) \\
Head noun min depth & $-0.0456$ & $0.0458$ & mixedlm (3/3) \\
Topological density & $0.0031$ & $0.0008$ & paired-t (2/3); mixedlm (1/3) \\
Scene-object ratio & $0.0047$ & $0.0033$ & mixedlm (3/3) \\
VG attr precision & $0.0357$ & $0.0049$ & mixedlm (3/3) \\
VG attr recall & $-0.0014$ & $0.0009$ & paired-t (3/3) \\
\bottomrule
\end{tabular}
\end{table}

The caption-length and projective-density coefficients are negative in all
three seeds, not just on average, so the direction of both headline effects is
consistent across seeds. This table is therefore not a sanity check appended to
the analysis: it is the only place the sign of each effect is established, and
the directional claims in the rest of this report rest on it.

A caveat on what Table~\ref{tab:mixed-effects} is actually doing for the
headline rows: as the Method column records, caption length falls back to the
per-image paired $t$-test in all three seeds, and projective spatial density
falls back in two of three. For those rows the point estimate equals the mean
of the per-image difference between the language-supervised and DINOv2
predictions and is unchanged by the fallback; what is \emph{not} done is the
image random-intercept absorption described in \S\ref{sec:mixed-effects}. The
direction claim still rests on a paired-by-image comparison, but not on a
fitted variance-components model. For the metrics that did fit
(\code{adj\_per\_noun}, \code{head\_noun\_min\_depth},
\code{scene\_object\_ratio}, \code{vg\_attr\_precision}, and topological
density in two seeds) the coefficients are from a true image random-intercept
fit.

\subsection{Effect Sizes Are Robust but Small}
\label{sec:magnitudes}

The effect size \(r\) measures how \emph{consistently} the per-image paired
difference takes one sign, not how \emph{large} that difference is. Each
encoder generates one caption for each of 5000 images per seed, so the
per-encoder sample size is 5000; but Wilcoxon's signed-rank test drops pairs
whose difference is exactly zero, so the effective $n$ that produces the
$r$ statistic is smaller. The ``Eff.\ $n$'' columns in
Tables~\ref{tab:beam-primary} and~\ref{tab:nucleus-primary} report this. For
caption length (a token count), most pairs have non-zero difference and the
effective $n$ is around 4000 per seed. For projective spatial density (a rate
that is often exactly zero), the effective $n$ is closer to 1100 per seed
under beam decoding. Both are still large samples, but the headline projective
result is supported by about one quarter as many image pairs as a reader
would assume from the 5000-per-encoder design. In native units, the two
robust beam effects are modest in absolute size.

\begin{table}[h]
\centering
\caption{Raw per-encoder means for the two robust beam metrics, averaged over
seeds 1--3. ``Lang.\ avg.'' averages CLIP and SigLIP. The final column is the
language-supervised minus DINOv2 difference and matches the beam coefficients
in the mixed-effects table.}
\label{tab:raw-means}
\begin{tabular}{@{}lrrrrr@{}}
\toprule
Metric & CLIP & SigLIP & DINOv2 & Lang.\ avg. & Lang.\ \(-\) DINOv2 \\
\midrule
Caption length (tokens) & 8.77 & 8.97 & 9.19 & 8.87 & \(-0.32\) \\
Projective spatial density & 0.0082 & 0.0107 & 0.0135 & 0.0094 & \(-0.0041\) \\
\bottomrule
\end{tabular}
\end{table}

DINOv2 captions average about one third of a token longer than the CLIP+SigLIP
average (9.19 vs.\ 8.87 tokens). On projective spatial density DINOv2 scores
0.0135 against a language-supervised average of 0.0094 --- small in absolute
terms, but a relative gap of roughly 44\%. ``Robust'' in this report therefore
means reproducible in direction across seeds and decoders, not practically
large: encoder pretraining reliably \emph{shifts} these caption properties
rather than transforming them.

\paragraph{Bootstrap uncertainty.} To put numeric uncertainty on the
per-seed $r$ statistics, we ran a 1000-iteration paired image-level
bootstrap (resampling per-image differences with replacement, recomputing
Wilcoxon $r$ at each iteration). The full table is in
Appendix~\ref{app:bootstrap}. The aggregate 95\% bootstrap CIs are
\textbf{[0.262, 0.324]} for beam projective density,
\textbf{[0.185, 0.218]} for beam caption length, and
\textbf{[0.132, 0.188]} for nucleus projective density --- all comfortably
above the pre-registered 0.10 floor. Two per-seed CIs reach the floor
honestly: beam caption length in seed 3 has CI $[0.099, 0.159]$, with the
lower bound a hair below 0.10, and nucleus projective density in seed 1
has CI $[0.077, 0.185]$. In both cases the aggregate is unambiguous and
the other two seeds are clearly above the floor, so the ``robust'' verdict
holds; but the per-seed near-misses are the bound on how confident a
reader should be from any single seed. Nucleus caption length, in
contrast, has aggregate CI $[0.026, 0.057]$ entirely below the floor ---
the ``not robust'' verdict is solidly below threshold, not borderline.

\subsection{Within-Category Structure}
\label{sec:within}

The primary contrast pools CLIP and SigLIP as a single ``language-supervised''
category. The pre-registered secondary family runs all three pairwise contrasts
(CLIP vs.\ DINOv2, SigLIP vs.\ DINOv2, CLIP vs.\ SigLIP) and is the right place
to check whether that pooling is innocuous. It is not entirely innocuous: CLIP
and SigLIP differ from DINOv2 at different magnitudes, and on one metric they
differ robustly from each other.

\begin{table}[h]
\centering
\caption{Pairwise contrasts on robust headline metrics and the metric on which
CLIP and SigLIP themselves differ robustly. Beam decoding, MAE excluded; mean
$r$ across seeds 1--3 (\code{captions/analysis\_aggregate\_beam\_no\_mae\_secondary.csv}).}
\label{tab:within}
\footnotesize
\begin{tabular}{@{}l rrr l@{}}
\toprule
Metric & CLIP$\leftrightarrow$DINOv2 & SigLIP$\leftrightarrow$DINOv2 & CLIP$\leftrightarrow$SigLIP & Note \\
\midrule
Projective spatial density & 0.358 \robust & 0.177 \robust & 0.188 \unstable & CLIP gap $\approx 2\times$ SigLIP gap \\
Caption length & 0.248 \robust & 0.139 \unstable & 0.121 \unstable & SigLIP--DINOv2 gap smaller \\
VG attr.\ recall & 0.207 \robust & 0.040 \notrobust & 0.256 \robust & CLIP differs from SigLIP \\
\bottomrule
\end{tabular}
\end{table}

Three observations follow.

\begin{itemize}
  \item On projective spatial density, both language-supervised encoders sit
  below DINOv2 in every seed, but CLIP is about twice as far from DINOv2 as
  SigLIP is ($\overline{r}=0.358$ vs.\ $0.177$). The pooled contrast is real;
  the within-category gap is not negligible.
  \item On caption length, the CLIP--DINOv2 gap is robust across seeds but the
  SigLIP--DINOv2 gap is only meaningful in two of three; the CLIP gap is about
  $1.8\times$ the SigLIP gap ($\overline{r}=0.248$ vs.\ $0.139$). CLIP carries
  most of the pooled length effect.
  \item On VG attribute recall, CLIP and SigLIP differ robustly from each
  other ($\overline{r}=0.256$, meaningful in 3/3 seeds), so pooling them
  conceals a real within-category difference. This metric is unstable in the
  primary contrast but stable in the pairwise contrast.
\end{itemize}

The cleanest reading of the headline is therefore narrower than ``the
language-supervised encoders differ from DINOv2.'' It is: \emph{CLIP differs
from DINOv2 strongly and consistently on both projective spatial language and
caption length; SigLIP differs from DINOv2 in the same direction but more
weakly, and on caption length the SigLIP--DINOv2 gap is not stable.} The
supervision-category framing is a useful first cut, but the per-pair
breakdown is what the data actually support.

\section{Interpretation}

The results support a precise version of the encoder-pretraining hypothesis:
visual encoder choice does not uniformly change every caption property, but it
does robustly affect specific linguistic signatures.

\begin{itemize}
  \item \textbf{Projective spatial language is robust, and DINOv2 leads.} The
  effect survives across three seeds and both decoding regimes in the
  MAE-excluded comparison; in every case the self-supervised encoder produces
  \emph{more} projective spatial language than the CLIP+SigLIP average. The
  pairwise breakdown (\S\ref{sec:within}) shows the gap is roughly twice as
  large for CLIP as for SigLIP, so the contrast is anchored at CLIP rather
  than uniform across the language-supervised category.
  \item \textbf{Caption length is robust under beam decoding only.} Under beam
  search DINOv2 captions are longer than the language-supervised average; the
  effect does not survive nucleus sampling, so it should be described as
  decoder-sensitive. Within the language-supervised category the
  CLIP--DINOv2 gap is robust across seeds but the SigLIP--DINOv2 gap is not
  (\S\ref{sec:within}), so the pooled length contrast is essentially a
  CLIP-vs-DINOv2 contrast.
  \item \textbf{Adjective density and noun specificity are not stable
  headlines.} They are meaningful in only one seed, so the final aggregate
  treats them as seed-sensitive.
  \item \textbf{MAE is a secondary result.} It repeatedly produces lower CIDEr
  under the adapter-only setup, so style differences involving MAE are
  confounded with quality.
\end{itemize}

The clean conclusion is:

\begin{quote}
Holding the decoder, adaptor architecture, training data, and decoding
procedure fixed, visual encoder pretraining changes downstream caption
language. The most robust signature is projective spatial language: the
self-supervised encoder (DINOv2) produces more of it than the
language-supervised encoders, across three seeds and both decoders, with the
gap roughly twice as large for CLIP as for SigLIP. Under beam search, DINOv2
also produces robustly longer captions; the per-pair breakdown shows that
length contrast is essentially CLIP versus DINOv2 (the SigLIP--DINOv2 gap is
not stable across seeds). Both differences are reliable in direction but
small in magnitude. MAE's lower CIDEr indicates that reconstruction-pretrained
features are less comparable in this adapter-only captioning setup and should
be treated separately from the primary comparable-quality claim.
\end{quote}

\section{Limitations}

\begin{enumerate}
  \item \textbf{Pretraining data is not controlled.} CLIP, SigLIP, DINOv2, and
  MAE differ in both objective and data source. The experiment compares real
  pretrained checkpoints, not models pretrained from scratch on matched data.
  Most notably, the two language-supervised encoders are pretrained on
  web-scraped image--text pairs while DINOv2 and MAE are pretrained on
  curated image collections; this web-versus-curated axis is confounded with
  the supervision-objective axis and is the most important residual confound
  on the headline contrast.
  \item \textbf{MAE quality confound.} MAE's lower CIDEr and SPICE mean it
  cannot support the same style-level claim as the comparable-quality
  encoders. Both reference metrics agree on this, so the MAE-exclusion
  decision is supported by two independent reference scorers rather than
  one.
  \item \textbf{Reference scoring depends on legacy Java.} SPICE 1.0 relies
  on the Nashorn JavaScript engine, which was removed from the JDK in
  version 15. The scoring run therefore requires JDK 11 (or earlier);
  attempting to run SPICE under a newer JDK produces \code{SPICE=n/a} and a
  null-pointer trace from Nashorn. This is a tooling fragility, not a
  correctness issue: the reported SPICE values use the upstream reference
  implementation under a compatible JDK. Reproducers should use JDK 11.
  \item \textbf{Frozen GPT-2 has its own priors.} GPT-2 medium may normalize
  outputs toward its pretrained language distribution, potentially muting
  encoder differences.
  \item \textbf{Small adaptor by design.} The lightweight adaptor preserves
  encoder differences, but a stronger Q-Former-like connector might improve
  quality or reduce observed differences.
  \item \textbf{The clean contrast has no self-supervised replication.}
  Excluding MAE leaves DINOv2 as the only self-supervised encoder in the
  headline comparison. The 2\(\times\)2 design was intended to provide
  within-category replication, so that a category-level effect could be
  distinguished from a single encoder's idiosyncrasies; in the
  comparable-quality contrast that safeguard is lost, and ``self-supervised''
  is confounded with ``DINOv2-specific.'' A second comparable-quality
  self-supervised encoder would be needed to restore it.
\end{enumerate}

\section{Recommended Reporting Language}

For a paper or presentation, the following wording is defensible:

\begin{quote}
We introduce a controlled frozen-decoder captioning probe for auditing how
visual encoder pretraining shapes generated language. With GPT-2, the adaptor,
COCO data, and decoding fixed, the comparable-quality contrast
(CLIP+SigLIP vs.\ DINOv2) shows a robust difference in projective spatial
language across three seeds and both decoding regimes: the self-supervised
encoder produces \emph{more} projective spatial language than the
language-supervised encoders, with the gap roughly twice as large for CLIP as
for SigLIP. Under beam decoding the same direction also holds for caption
length, but the per-pair breakdown shows that length contrast is essentially
a CLIP-vs-DINOv2 effect. Both effects are consistent in direction and above
our pre-registered effect-size floor in every seed, though small in absolute
size. MAE produces lower CIDEr across seeds under the adapter-only setup, so
we report it as a secondary finding rather than part of the headline
comparable-quality contrast.
\end{quote}

\appendix

\section{Reproducibility Pointers}

\begin{table}[h]
\centering
\caption{Key scripts and result files.}
\label{tab:repro}
\begin{tabularx}{\textwidth}{@{}l X@{}}
\toprule
Path & Role \\
\midrule
\code{scripts/02\_train\_clip.py} & Trains one adaptor while freezing the selected visual encoder and GPT-2. The filename retains an early CLIP-only label; the script is the general adaptor trainer and is invoked once per encoder (CLIP, SigLIP, DINOv2, MAE). \\
\code{scripts/04\_generate\_captions.py} & Generates beam and nucleus captions for a selected encoder and seed. \\
\code{scripts/05\_score\_metrics.py} & Computes per-caption style and attribute metrics. \\
\code{scripts/06\_analyze.py} & Runs per-seed primary, secondary, and image-controlled analyses. \\
\code{scripts/08\_caption\_quality.py} & Computes the CIDEr and SPICE quality preconditions. Requires JDK 11 because SPICE 1.0 uses Nashorn. \\
\code{scripts/09\_aggregate\_seeds.py} & Aggregates per-seed analyses into cross-seed robustness summaries. \\
\code{scripts/10\_bootstrap\_ci.py} & 1000-iteration paired image-level bootstrap of the Wilcoxon $r$ statistic; produces \code{captions/bootstrap\_ci.csv}. \\
\code{captions/analysis\_aggregate\_beam\_no\_mae\_primary.csv} & Cleanest headline beam result. \\
\code{captions/analysis\_aggregate\_nucleus\_no\_mae\_primary.csv} & Cleanest headline nucleus result. \\
\bottomrule
\end{tabularx}
\end{table}

\section{Per-Seed Statistical Detail for Headline Metrics}
\label{app:bootstrap}

Table~\ref{tab:per-seed-detail} reports, for each headline metric and
decoder, the Benjamini--Hochberg adjusted $p$-value, the Wilcoxon effective
sample size, the point estimate of $r$, and the 95\% bootstrap CI on $r$.
The bootstrap is paired-by-image with 1000 iterations; the CI is the
percentile method. Per-seed rows show within-seed statistics; the
``aggregate'' row is the cross-seed mean $r$ with its corresponding
bootstrap CI (independently resampling each seed and averaging the three
$r$ values per iteration). No BH correction is applied across seeds, so
the aggregate row has no $p_{\text{adj}}$.

\begin{table}[h]
\centering
\caption{Per-seed BH-adjusted $p$-values and bootstrap 95\% CIs for the
robust headline metrics (and the nucleus caption-length row that is not
robust, for completeness). Sources:
\code{captions/analysis\_seed*\_*\_no\_mae\_primary.csv} (BH-adjusted $p$
and effective $n$) and \code{captions/bootstrap\_ci.csv} (bootstrap CIs).}
\label{tab:per-seed-detail}
\footnotesize
\begin{tabular}{@{}l l l r r r l@{}}
\toprule
Decoder & Metric & Seed & Eff.\ $n$ & $r$ & 95\% CI & $p_{\text{adj}}$ \\
\midrule
beam & projective density & 1 & 1027 & 0.276 & [0.221, 0.332] & $3.7\!\times\!10^{-18}$ \\
beam & projective density & 2 & 1134 & 0.256 & [0.198, 0.314] & $3.0\!\times\!10^{-17}$ \\
beam & projective density & 3 & 1154 & 0.345 & [0.295, 0.396] & $8.3\!\times\!10^{-31}$ \\
beam & projective density & agg. & --- & 0.292 & [0.262, 0.324] & --- \\
\midrule
beam & caption length & 1 & 4103 & 0.191 & [0.162, 0.222] & $1.4\!\times\!10^{-33}$ \\
beam & caption length & 2 & 4096 & 0.287 & [0.261, 0.316] & $2.0\!\times\!10^{-74}$ \\
beam & caption length & 3 & 4027 & 0.128 & [0.099, 0.159] & $1.8\!\times\!10^{-15}$ \\
beam & caption length & agg. & --- & 0.202 & [0.185, 0.218] & --- \\
\midrule
nucleus & projective density & 1 & 1570 & 0.131 & [0.077, 0.185] & $1.8\!\times\!10^{-6}$ \\
nucleus & projective density & 2 & 1600 & 0.175 & [0.126, 0.226] & $2.2\!\times\!10^{-11}$ \\
nucleus & projective density & 3 & 1641 & 0.175 & [0.126, 0.222] & $9.8\!\times\!10^{-12}$ \\
nucleus & projective density & agg. & --- & 0.160 & [0.132, 0.188] & --- \\
\midrule
nucleus & caption length & 1 & 4685 & 0.057 & [0.029, 0.086] & $2.3\!\times\!10^{-4}$ \\
nucleus & caption length & 2 & 4688 & 0.005 & [0.001, 0.031] & 0.80 \\
nucleus & caption length & 3 & 4699 & 0.051 & [0.022, 0.080] & $1.8\!\times\!10^{-3}$ \\
nucleus & caption length & agg. & --- & 0.038 & [0.026, 0.057] & --- \\
\bottomrule
\end{tabular}
\end{table}

Two reading notes. (i)~Per-seed CIs that include or touch the 0.10
pre-registered floor are honest near-misses: beam caption length in
seed~3 has lower bound 0.099; nucleus projective density in seed~1 has
lower bound 0.077. In both cases the aggregate CI is cleanly above 0.10
and the other two seeds clear the floor by a wide margin. (ii)~The
nucleus caption-length row is included because, although the
BH-adjusted $p$-values are significant in two of three seeds, the effect
sizes are an order of magnitude below the floor and the aggregate CI
$[0.026, 0.057]$ is entirely below it --- which is exactly the case the
pre-registered effect-size floor was designed to filter out.

\begin{thebibliography}{9}

\bibitem{radford2021clip}
Alec Radford et al.
\newblock Learning transferable visual models from natural language supervision.
\newblock \emph{ICML}, 2021.

\bibitem{zhai2023siglip}
Xiaohua Zhai et al.
\newblock Sigmoid loss for language image pre-training.
\newblock \emph{ICCV}, 2023.

\bibitem{oquab2023dinov2}
Maxime Oquab et al.
\newblock DINOv2: Learning robust visual features without supervision.
\newblock arXiv:2304.07193, 2023.

\bibitem{he2022mae}
Kaiming He et al.
\newblock Masked autoencoders are scalable vision learners.
\newblock \emph{CVPR}, 2022.

\bibitem{vedantam2015cider}
Ramakrishna Vedantam, C. Lawrence Zitnick, and Devi Parikh.
\newblock CIDEr: Consensus-based image description evaluation.
\newblock \emph{CVPR}, 2015.

\bibitem{anderson2016spice}
Peter Anderson et al.
\newblock SPICE: Semantic propositional image caption evaluation.
\newblock \emph{ECCV}, 2016.

\end{thebibliography}

\end{document}
